% easychair.tex,v 3.5 2017/03/15

\documentclass{easychair}
%\documentclass[EPiC]{easychair}
%\documentclass[EPiCempty]{easychair}
%\documentclass[debug]{easychair}
%\documentclass[verbose]{easychair}
%\documentclass[notimes]{easychair}
%\documentclass[withtimes]{easychair}
%\documentclass[a4paper]{easychair}
%\documentclass[letterpaper]{easychair}

\usepackage{doc}
\usepackage{minted}
%% \usepackage{amsfonts}
\usepackage{graphicx}
\usepackage{amsmath}
\usepackage{amssymb}
\usepackage{bussproofs}
\usepackage{cite}

% For ≞ (requires the LuaLaTeX engine)
\usepackage{unicode-math}
\setmathfont{Stix Two Math}

% Use this if you have a long article and want to create an index
% \usepackage{makeidx}

% In order to save space or manage large tables or figures in a
% landcape-like text, you can use the rotating and pdflscape
% packages. Uncomment the desired from the below.
%
% \usepackage{rotating}
% \usepackage{pdflscape}

% Some of our commands for this guide.
%
\newcommand{\easychair}{\textsf{easychair}}
\newcommand{\miktex}{MiK{\TeX}}
\newcommand{\texniccenter}{{\TeX}nicCenter}
\newcommand{\makefile}{\texttt{Makefile}}
\newcommand{\latexeditor}{LEd}
\newcommand{\Theory}{\texttt{Theory}}
\newcommand{\Proof}{\texttt{Proof}}
\newcommand{\Proposition}{\texttt{Proposition}}
\newcommand{\Bool}{\texttt{Bool}}
\newcommand{\arrow}{\to}
\newcommand{\limp}{\Rightarrow}
\newcommand{\True}{\texttt{True}}
\newcommand{\False}{\texttt{False}}
\newcommand{\STV}[2]{<\!#1, #2\!>}
\newcommand{\arrowtype}{\text{->}}
\newcommand{\axon}{\text{ax-1}}
\newcommand{\axtw}{\text{ax-2}}
\newcommand{\axth}{\text{ax-3}}
\newcommand{\axmp}{\text{ax-mp}}

%\makeindex

%% Front Matter
%%
% Regular title as in the article class.
%
%% \title{Estimating the Probability of a Proposition to be Provable
%%   Using Probabilistic Logic Networks}
%% \title{Estimating the Probability of a Conjecture to be a Theorem with
%%   Probabilistic Logic Networks, and How to Take Advantage of it for
%%   Inference Control}
\title{Using Forward Chaining to Go Backward}

% Authors are joined by \and. Their affiliations are given by \inst, which indexes
% into the list defined using \institute
%
\author{Nil Geisweiller}

% Institutes for affiliations are also joined by \and,
\institute{
  SingularityNET Foundation,\\
  Zug, Switzerland\\
  \email{nil@singularitynet.io}
}

%  \authorrunning{} has to be set for the shorter version of the authors' names;
% otherwise a warning will be rendered in the running heads. When processed by
% EasyChair, this command is mandatory: a document without \authorrunning
% will be rejected by EasyChair
\authorrunning{Geisweiller}

% \titlerunning{} has to be set to either the main title or its shorter
% version for the running heads. When processed by
% EasyChair, this command is mandatory: a document without \titlerunning
% will be rejected by EasyChair

\titlerunning{Using Forward Chaining to Go Backward}

\begin{document}

\maketitle

%% \begin{abstract}
%% \end{abstract}

\section{Introduction}

If all you have is a forward chainer, is it be possible to emulate a
process that is functionally equivalent to backward chaining?  This is
the question we attempt to answer in this paper.  Such question stems
from the observation that some rewriting systems, such as
MORK~\cite{NEXT}, perform better on forward chaining than on backward
chaining~\cite{NEXT}.  This is to be contrasted with some other
rewriting systems such as PeTTa~\cite{NEXT}, a Prolog-based rewriting
system, which performs equally well on forward and backward
chaining~\cite{NEXT}.  Since forward chaining in MORK outperforms
forward chaining in PeTTa~\cite{NEXT}, it gives us an incentive to
emulate backward chaining using forward chaining to obtain a
MORK-based backward chainer that outperforms PeTTa.

\section{Forward vs Backward Chaining}

Forward chaining is a reasoning process that begins with axioms and
apply inference rules to produce more truths.  For instance, using
axioms and inference rules of propositional calculus (borrowed from
Metamath~\cite{NEXT}), given with a type theoretic representation
\begin{itemize}
\item $\axon : \phi \to (\psi \to \phi)$ $\ \ \ \ \ \ \ \ \ \ \ \ \ \
  \ \ \ \ \ \ \ \ $ $\axtw : ( \phi \to ( \psi \to \chi ) ) \to ( ( \phi \to \psi ) \to ( \phi \to \chi ) )$
\item $\axmp : \phi\ \arrowtype\ ( \phi \to \psi )\ \arrowtype\ \psi$
\end{itemize}
where $\arrowtype$ represents the arrow type, one can build a proof of
the identity in a forward way:
\begin{enumerate}
\item apply $\axmp$ to $\axon$ and $\axtw$ to obtain
  $(\axmp\ \axon\ \axtw) : ( \phi \to \psi) \to (\phi \to \phi)$,
\item apply $\axmp$ to $\axon$ and step 1 to obtain
  $(\axmp\ \axon\ (\axmp\ \axon\ \axtw)) : \phi \to \phi$.
\end{enumerate}
On the contrary, a backward chainer starts with the proposition to be
proven, and expands the proof backward until reaching the axioms.  For
instance, the process may start with $x? : \phi \to \phi$ where $x?$
is a hole, then expand $x?$ backward with $\axmp$ to obtain
$(\axmp\ y?\ z?) : \phi \to \phi$, and keep expanding until reaching
$(\axmp\ \axon\ (\axmp\ \axon\ \axtw)) : \phi \to \phi$.  The
advantage of backward chaining is that the search space is constrained
by proofs that can only prove the target proposition, resulting in
considerable pruning.  The inconvenient is that it is more
sophisticated, involving full syntactic unification and maintaining
hypothetical paths.

\section{Emulating Backward Chaining with Forward Chaining}

To emulate backward chaining using forward chaining NEXT.

To achieve this feat we decompose the backward chaining task into
three steps
\begin{enumerate}
\item Use forward chaining to generate a theory of backward expansion
  up to some depth.
\item Given some target theorem, apply forward chaining on the theory
  of backward expansion to generate composite rules, such that, if
  they were to run forward, they would derive the given theorem only.
\item Use forward chaining on the obtained composite rules to generate
  proofs of the target theorem.
\end{enumerate}
You may have noticed that these three steps only make use of forward
chaining.

\newpage

%------------------------------------------------------------------------------

\bibliographystyle{plain}
\label{sect:bib}
%\bibliographystyle{alpha}
%\bibliographystyle{unsrt}
%\bibliographystyle{abbrv}
\bibliography{BackwardViaForward}

\newpage

%------------------------------------------------------------------------------

\appendix

%------------------------------------------------------------------------------
\end{document}
